\section{Limitations \& Future Prospects}

Although this paper has provided initial evidence for the feasibility, competitiveness, and promising scaling potential of \method for text generation in continuous latent space, we view it as a starting point for further exploration rather than a finished endpoint. First, at the scale and evaluation level, the current results reveal encouraging trends, but the experiments are still conducted at a relatively controlled scale and mainly serve to clarify the key properties of the framework. It is therefore natural and important to further examine its upper bound under larger model sizes, longer training, and more substantial compute budgets. Second, at the model-design level, our analyses show that the training strategy of the Text VAE, the text compression scheme, the choice of latent dimensionality, the semantic smoothness of the latent space, and the joint calibration of VAE logSNR, DiT block size, and noise schedule all affect the semantic organization of the latent space and the final generation quality. In particular, the experiments suggest that stronger latent representations usually require better-aligned noise calibration, indicating substantial room for further optimization. Finally, at the framework level, the main value of \method lies not merely in the denoising process itself, but in its decomposition of text generation into global semantic prior modeling and local textual realization. This opens the door to exploring stronger latent modules, such as AE \citep{bank2023autoencoders} and RAE \citep{zheng2025diffusion}, as well as more flexible prior-learning approaches, such as drifting-model-based \citep{deng2026generative} distribution matching for continuous priors. More broadly, following the idea of unified continuous latent-space modeling, the framework may also be extended to continuous modalities such as images, further advancing unified generation.

\section{Conclusion}

In conclusion, this paper presents \method, a hierarchical continuous latent diffusion language model that decomposes text generation into global semantic prior modeling in latent space and local textual realization through conditional decoding, thereby providing a principled alternative to strictly token-level language modeling. Across the full study, both the theoretical analysis and the experiments consistently suggest that text generation can benefit from hierarchical information decomposition: we find evidence of shared global semantic structure in latent space, identify effective design choices for latent-space formation and diffusion modeling, and show that under strictly matched comparisons, \method exhibits strong generation quality and encouraging scaling behavior. More broadly, our results indicate that for this class of models, generation-oriented evaluation and scaling trends may be more informative than likelihood alone, while the continuous latent formulation also offers a concrete path toward more native unified modeling across discrete text and continuous modalities.